\section{Numerical details of lattice study.}

In the previous section we described what our strategy was to
address the question of existence of a continuous limit for the
operators and the large $N$ transition point. Also, we made an
effort to convince ourselves that the large $N$ transitions indeed
are continuous. We now proceed to describe the exact details of
our simulations and the numbers that came out.

\subsection{The determination of $f_c(l)$.}

\subsubsection{Estimating the  gap at infinite $N$.}

At infinite $N$, small $\hat W$ loops will have a spectral gap on
the lattice. We need a method to extract that gap from numerical
data at finite $N$. As we dilate the loop, and/or reduce the
amount of smearing, the gap shrinks. We need a way to extrapolate
to the point where the gap vanishes. We wish to keep away from
measurements for too small gaps, and enter that regime only by
extrapolation.

Getting the infinite $N$ value of the gap is difficult because
subleading effects are large at the gap's edge. It is reasonable
to assume that the universal behavior at the gap's edge is the
same as at the edge of a gap in simple hermitian matrix models.
Under this assumption, we have separate predictions for the
behavior of the eigenvalue density and for the extremal
eigenvalue. One can check numerically how well the universal forms
work and we find them to work quite well. To fit onto the
universal behavior we need to adjust a parameter that corresponds
to the gap size at infinite $N$. This provides the determination
of the gap sizes we use subsequently.

It is only on the side of the transition where the eigenvalue density
has a gap that
there is a regime where large $N$ corrections are dominated by
universal functions; therefore, we are able to get good estimates
for the eigenvalue gap size at infinite $N$ from data obtained at one single,
but large $N$. For good measure, we checked that the predicted rational
power dependence on $N$ holds in the regimes in which we did our
analysis. A comparison with some simulations at $N=28$ convinced
us that the power dependence on $N$ is already at its known
universal value at
$N=28,44$.

We studied the eigenvalue
distribution of $\hat W [L , L ; f;
n=L^2;N]$ on the lattice
for several values of $L$ and $b$ at a fixed large $N$.
Let $e^{i\theta_k}, k=1,\cdots,N$ be the $N$ eigenvalues
of a fixed $L\times L$ loop on the lattice, $-\pi \le \theta_1 <
\theta_2 < \cdots < \theta_{N-1} < \theta_N < \pi$.
The sum of all angles is zero mod $2\pi$. On a lattice of $V$
sites there are $6V$ different $L\times L$ Wilson loops with
a fixed orientation; we computed the eigenvalues
of all such loops. In this manner, we obtain a
histogram of the distribution
$p_f(\theta;U)$ of all eigenvalues
in a fixed gauge field background $U$.
We also collected data for a histogram of the distribution
$p_m(\theta_N;U)$ for the largest eigenvalue in a fixed
gauge field background. In addition, we record
the mean, variance, skewness and kurtosis
pertaining to the largest eigenvalue:
\begin{eqnarray}
\mu_N(U) &=& \langle \theta_N \rangle_U;\nonumber\\
\sigma^2_N(U) &=& \langle \theta^2_N \rangle_U - \langle
\theta_N \rangle^2_U;\nonumber\\[0.6em]
S_N(U) &=&
\frac{ \langle \theta^3_N \rangle_U
-3 \langle \theta^2_N \rangle_U \langle \theta_N \rangle_U
+ 2 \langle \theta_N \rangle^3_U }{\sigma^3_N(U)};\nonumber\\[0.6em]
K_N(U) &=&
\frac{ \langle \theta^4_N \rangle_U
-4 \langle \theta^3_N \rangle_U \langle \theta_N \rangle_U
+6 \langle \theta^2_N \rangle_U \langle \theta_N \rangle^2_U
-3 \langle \theta_N \rangle^4_U }{\sigma^4_N(U)}-3.
\end{eqnarray}

If the eigenvalue distribution of $\hat W$ has a sizable gap, it
is reasonable to ignore the influence of the two edges on each
other. Then, the universal behavior at the gap's edge ought to be
the same as that of the edge of the gap in simple Gaussian
hermitian matrix model. Under this assumption, we compared our
data with the behavior of the extremal eigenvalue~\cite{tracy}
and the eigenvalue density close to the edge~\cite{gff}.

The universal distribution of the
extremal eigenvalue, $p(s)$, is given by~\cite{tracy}
\be
p(s)=\left [ \int_s^\infty q^2(x) dx \right ] e^{-\int_s^\infty (x-s)q^2(x) dx}
\ee
with $q$ being the solution to the Painleve II equation,
\be
\frac{d^2q}{ds^2} = sq + 2 q^3;
\ee
satisfying the asymptotic condition
\be
q(s) \sim Ai(s)\ \ {\rm as}\ \ s\rightarrow\infty.
\ee
The mean, variance, skewness and kurtosis of the above
distribution are $-1.7710868074$, $0.8131947928$,
$0.2240842036$ and $0.0934480876$. In order to match
the lattice data for the distribution of $\theta_N$
to the above universal form,
we define $s$ according to
\be
\theta_N = \frac{s}{\alpha} + \pi(1-g)
\label{scalevar}
\ee
and set
\begin{eqnarray}
\alpha &=& \sqrt{\frac{0.8131947928}{\sigma^2_N}}\nonumber\\
g &=& 1-\frac{1}{\pi}[\mu_N + \frac{1.7710868074}{\alpha}]
\end{eqnarray}
$\alpha$ is a non-universal scale factor and $g$ is half the arc
length of the gap on in the eigenvalue distribution on the unit
circle, centered at $-1$. The quantities, $\sigma^2_N$ and $\mu_N$
are the averages over the gauge field ensemble of the previously
defined quantities, $\sigma^2_N(U)$ and $\mu_N(U)$. From the
lattice data for $\sigma^2_N(U)$ and $\mu_N(U)$, we obtain an
estimate for $\alpha$ and $g$ with errors extracted by the
jackknife method. Using these estimates, we compared the scaled
lattice data with $p(s)$ for fixed $f$, $b$, $N$ and $V$. One such
comparison, at $f=0.15$, $b=0.361653$, $L=4$, $N=44$ and $V=7^4$
with a gap of $0.1229(4)$, is shown in Figure~\ref{gap} and there
is good agreement. In spite of the fact that data is needed over
the entire range of $s$ to get a good estimate of the skewness and
kurtosis, we find agreement between lattice data and the universal
estimates to within $5\%$ for the skewness and to within $20\%$
for the kurtosis for a wide range of gap values. The estimate of
the skewness and kurtosis do not depend on the extracted values of
$\alpha$ and $g$.

\begin{figure}[ht]
\centering
\includegraphics[width=\textwidth]{images/gap.pdf}
\caption{ A comparison of the
distribution of scaled extremal eigenvalue obtained on the lattice
(circles) at $f=0.15$, $b=0.361653$, $L=4$, $N=44$ and $V=7^4$
with that of the universal distribution (solid curve). }
\label{gap}
\end{figure}

The main result of the above numerical analysis is the
half gap size $g$, for fixed $f, L, b$, $g(f,L,b)$.


There is another way to obtain an estimate for $g$:
A recent result
provides the asymptotic corrections to the eigenvalue
density of Gaussian hermitian matrices~\cite{gff}. Using that
formula, we fitted
he eigenvalue density in terms of the scaled variable
in (\ref{scalevar}) to
\begin{equation}
p_f(s) = c \left\{ [Ai^\prime(s)]^2 - s [Ai(s)]^2 \right\} +d
\left\{ 3s^2[Ai(s)]^2 - 2s [Ai^\prime(s)]^2 -3Ai(s)Ai^\prime(s)
\right\} \label{airy}
\end{equation}
The variable $c$ is fixed by the overall
normalization and $\frac{d}{c}=-\frac{1}{20N^{2/3}}$ in the
asymptotic formula~\cite{gff} where $N$ is the size of the matrix.
The parameter $d$ could be non-universal; also $c$ could in some
sense be non-universal, since it depends on the total number of
eigenvalues, and it is not clear whether the Gaussian hermitian
matrix model matches onto our model with the same $N$ in both
cases. Therefore, we fitted the lattice distribution to the above
asymptotic formula with $c$ and $d$ as independent variables in
addition to the two variables appearing in (\ref{scalevar}),
producing a second determination of $g(f,L,b)$.

\begin{figure}[ht]
\centering
\includegraphics[width=\textwidth]{images/dist.pdf}
\caption{A comparison of the distribution of eigenvalue distribution
obtained on the lattice (circles with dashed lines) at $f=0.15$,
$b=0.361653$, $L=4$, $N=44$ and $V=7^4$ with that of the universal
distribution near the edge (solid curve). }
\label{dist}
\end{figure}

Figure~\ref{dist} shows the comparison of the full eigenvalues
distribution with the universal distribution near the edge. The
universal distribution captures the tail for the same reason that
the distribution of the extremal eigenvalue is in agreement with
$p(s)$. But we also see good agreement with the first two
oscillations near the edge and this includes more than the
extremal eigenvalue. In general we found the determination of the
gap to be consistent with the one obtained using only the extremal
eigenvalue. For an error estimate, we only used the extremal
eigenvalue method.



\subsubsection{Extrapolating to zero gap.}



The next step is to extract the critical value of $f$ at a
fixed $b$ and $L$ at infinite $N$. This is done by
assuming a fit of the form
\be
g(f,L,b) = A(L,b) \sqrt{f - F_c(L,b)}
\ee
The square root dependence on external parameters that enter the
probability law in an analytic manner is generic in matrix models
and well supported by the data.

A plot of $g^2(f,L,b)$ as function a $f$ at $b=0.361653$, $N=44$
is shown in figure~\ref{gap_vs_ff}. $F_c(L,b)$ is extracted using
a range of $f$ such that the gap, $g(f,L,b)$, is not too close to
$0$ nor is it too far from $0$. If we go to too small gaps we fear
unaccounted for rounding effects; if the gap is too large, there
is little reason to use the universal square root behavior we
postulated.

Figure~\ref{gap_vs_ff} also shows the fit along with the estimate
for $F_c(L,b)$.


\begin{figure}[ht]
\centering
\includegraphics[width=\textwidth]{images/gap_vs_ff.pdf}
\caption{ Plot of
$g^2(f,L,b)$ as a function of $f$ at $b=0.361653$, $N=44$ and
$V=7^4$. A fit to the data along with the extrapolated value for
$F_c(L,b)$ is also shown. }
\label{gap_vs_ff}
\end{figure}


\subsubsection{The continuum limit of $F_c(L,b)$.}


We now proceed to take the continuum limit of $F_c(L,b)$.
We used two loop tadpole improved perturbation theory
to keep the physical loop fixed as we varied $b$.
Using~\cite{knn}, we have a numerical formula for the critical
deconfinement temperature in lattice units, in terms of tadpole
improved perturbation theory, with $e(b)$ being the average value of the
plaquette for untwisted single plaquette gauge action at coupling
$b$.
\begin{equation}
\frac{1}{T_c(b)} = 0.26
\left[\frac{11}{48\pi^2 b_I}\right]^{\frac{51}{121}}
e^{\frac{24\pi^2 b_I}{11}};\ \ \ \
b_I=be(b)
\label{twoloop}
\end{equation}
This will be used to determine the lattice spacing in physical units.

We set $lt_c=LT_c(b)$ and accumulate data at several values of the
physical loop size: $lt_c=0.42,0.48,0.58,0.64$. At each of these
physical sizes, we went through a sequence of three to four
lattice spacings $a(b)$, and, for each $b$, we obtained a value
for $F_c(L=\frac{l}{a(b)},b)$, using the methods explained above.

For each $lt_c$, the $F_c(L=\frac{l}{a(b)},b)$ values were then
extrapolated to the continuum assuming
\be
F_c(L=\frac{l}{a(b)},b) = f_c(lt_c) + O(T_c^2(b)).
\ee
where we recall
that $T_c(b)=a(b)t_c$. (We shall use the slightly inconsistent
notations $f_c(l)$ and $f_c(lt_c)$ for the same quantity, where the
argument is either dimensional or made dimensionless by the appropriate
power of the physical deconfining temperature $t_c$.)

The extrapolated results for $f_c(lt_c)$ are
shown in figure~\ref{crit2}.


\begin{figure}[ht]
\centering
\includegraphics[width=\textwidth]{images/crit2.pdf}
\caption{ Plot of $F_c(L=\frac{l}{a(b)},b)$ as a function of $T_c(b)$ for
four different values of $lt_c$. Also shown are fits to the data
producing continuum values for $f_c(lt_c)$. }
\label{crit2}
\end{figure}



Our final results are plotted in the $(f,l)$ plane. The critical
line $f_c(l)\equiv f_c (lt_c)$ divides this plane into two as
shown in figure~\ref{fltc}. In the upper part Wilson loops have a
gap in their eigenvalue distribution, whereas in the lower part
there is no gap. The band bounded from below
by the critical zero-lattice-spacing line indicates the range
of finite-lattice-spacing effects.



\subsection{A study of both sides of the transition under scaling.}

Finite lattice size effects are evident in figure~\ref{fltc}. It
is impractical to just scale the loop because the curve $f_c(l)$
varies when $l$ is changed in the appropriate range by less than
the lattice corrections. To see the effect of scaling we need to
go on a line in the $(f,l)$ plane that is not horizontal, which
means we are going to vary the smearing as a function of the
physical loop size. As usual, we trade $l$ for the dimensionless
variable $lt_c$.

We chose the line (we shall refer to this line as the ``crossing line'')
\begin{equation}
f = 0.25 - (0.6128 lt_c)^2
\label{fline}
\end{equation}
shown in figure~\ref{fltc} as it cuts across the entire finite band of lattice
corrections bounded by the critical line $f_c(lt_c)$.


Consider a fixed loop of size $L$ on the lattice for various
lattice couplings $b$ such that $f(b)$ is given by (\ref{fline})
with $lt_c=LT_c(b)$ and $T_c(b)$ given by (\ref{twoloop}). We wish
to study the nature of the transition in the eigenvalue
distribution along this line.

\begin{figure}[ht]
\centering
\includegraphics[width=\textwidth]{images/fltc.pdf}
\caption{ The critical line in the
$(f,lt_c)$ plane. Also shown is the crossing line used to
investigate the phase transition induced by scaling the Wilson
loops. }
\label{fltc}
\end{figure}

We already know that the distribution should have no gap for $l
> l_c$ whereas it should have a gap for $l < l_c $. In
addition to getting an estimate for $l_ct_c$, we can also
investigate the universality class of this transition by keeping
$L$ fixed and going along the crossing line given by
(\ref{fline}). We set $L=3,4,5$ and worked in the following ranges
of $b$: for L=3, in $0.345 < b < 0.365$, for L=4, in $0.346 < b <
0.368$ and for L=5, in $0.351 < b < 0.369$. In each range, as $b$
decreases the physical loop is scaled up. In this way one gets
overlaps at the same physical loop size, once represented by a
smaller value of $L$ and a smaller value of $b$, and other times,
by correspondingly larger values of these parameters. Because of
scaling violations one will not get identical eigenvalue
distributions; we need to see that, as $b$ and $L$ increase, these
corrections decrease and that a finite limiting eigenvalue
distribution is approached. This is the eigenvalue distribution
associated with the continuum loop of size $l$.

In Figure~\ref{fltc} we display, for each value of $L=3,4,5$, the
two critical lengths one obtains when extrapolating from each
phase independently as triangles on the crossing line. As $L$ is
increased, we see that the distance between the two members of
each pair decreases as $L$ increases. Also, the middle points of
the segments connecting the points in each pair appear to converge
as $L$ increases. This indicates that in the continuum limit there
will be one single transition point on this line, a loop size at
which a continuous transition in the large $N$ spectrum of ${\hat
W}$ occurs.

\subsubsection{Estimating the critical size from the small loop
side.}\label{gapside}

We again obtained estimates for the infinite $N$ gap, using the
extremal eigenvalue distribution ($p(s)$) and the eigenvalue density
($p_f(s)$) near
the edge. The estimates for the gap came out consistent with each other
and the results are plotted in figure~\ref{gap_diag}. The
critical size of the loop is extracted by again assuming
a square root singularity written in the form
\be
g^2(lt_c)= A \ln \frac{l}{l_c} \ee
For each value of $L=3,4,5$ we
got an estimate for $l_ct_c$ which is shown in
figure~\ref{gap_diag}. One sees that the lines $L=4,5$ are closer
to each other than the lines $L=3,4$; the ratio of the approximate
distances between the two pairs of lines is in rough agreement
with a hypothesis of order $a^2 (b)$ corrections.

\begin{figure}[ht]
\centering
\includegraphics[width=\textwidth]{images/gap_diag.pdf}
\caption{ The square of the
gap as a function of $\ln(lt_c)$ along the crossing line. }
\label{gap_diag}
\end{figure}


\subsubsection{Estimating the critical size from the large loop
side.}\label{rhopi}



For operators with no gap we need a method to extract the infinite
$N$ value of the eigenvalue density at angle $\pi$, $\hat\rho
(\pi)$. Shrinking the loop and/or smearing it more makes $\hat\rho
(\pi)$ decrease. We also need a method to extrapolate to the point
where $\hat\rho(\pi)$ vanishes, without getting too close to it.

The estimation of the infinite $N$ value of $\hat \rho(\pi)$ also
suffers from large subleading effects, of order $\frac{1}{N}$
typically. There is a rapid fluctuation in the eigenvalue density
reflective of the repulsion between eigenvalues. The rapid
fluctuation is eliminated by Fourier transforming the numerically
obtained binned eigenvalue distribution with respect to the angle
for wave-lengths kept large relative to $\frac{1}{N}$ . This is
equivalent to evaluating the traces of $\hat W^k$ for consecutive
values of $k$, all kept much smaller than $N$. Keeping only those
Fourier modes, we get our estimates for $\hat \rho(\pi)$; checks
against a simulation at lower $N$'s shows that we are determining
$\hat \rho(\pi)$ to reasonable accuracy.


The extrapolation of $\hat \rho(\pi)$ to zero as the loop size is
decreased and/or its amount of smearing is increased is done in a
similar way as when we deal with the gap. We find that the leading
behavior is again square-root like, but this holds in a small
range which can be extended by adding more powers of $(l-l_c)$.
The determination of $l_c$ by this method is seen to be quite
robust.

In two dimensions the large $N$ eigenvalue distribution for
loops of arbitrary size is known in the continuum \cite{duol}.
As we shall elaborate later, it is plausible to assume
that the behavior close to the large $N$
transition is universal, and holds also for
our $\hat \rho(\theta)$ in four dimensions.
Therefore, we
postulate that
${\hat \rho(\pi)}$
vanishes as a square root
at the critical value of any parameter
that enters the
underlying distribution controlling
the smeared Wilson loop operators in an
analytic manner. This is why we took $\hat\rho (\pi)$ to vanish as
$\sqrt{l-l_c}$ as one approaches
the appropriate critical size, $l_c$, from above.


As indicated above, for $l$ large enough with respect to $l_c$
we extract the nonzero value of $\hat\rho(\pi)$ by first
representing the distribution for all angles as:
\be
\hat\rho(\theta) = \sum_{k=0}^N f_k \cos(k\theta)
\ee
Next, we truncate the Fourier series and use the truncated series
to provide the value of $\hat\rho(\pi)$.
\be
\hat\rho(\pi) = \sum_{k=0}^{k_m} (-1)^k f_k
\ee
The value of $k_m$ is
chosen to be slightly less than $N/2$; this eliminates
oscillations of frequency $N$. We checked that our choice for
$k_m$ was reasonable by looking for stability as we vary $k_m$.

We
found that for low enough $k$'s $f_k$ oscillates in sign from
$k=1$ onwards.
As Figure~\ref{fourier} shows, for $k$ up to about $\frac{N}{2}$
the coefficients are smooth in $k$, and the sign oscillation is
obeyed. We shall come back to this property later.

\begin{figure}[ht]
\centering
\includegraphics[width=\textwidth]{images/fourier.pdf}
\caption{ The fourier
coefficients $f_k$ as a function of $k$ for two different loops in
the gapless region. The measurements were obtained using an $L=5$
loop at $N=44$ at two different lattice spacings. }
\label{fourier}
\end{figure}



The values of $\hat\rho(\pi)$ we obtain are fitted to a square root
in the critical region and the results are shown in figure~\ref{rho_diag}.
The estimates for $l_ct_c$ from $\hat\rho(\pi)$ tend to be higher
than the ones coming from fitting the gap. This difference
decreases as the spacing goes to zero.

The critical values for $l_ct_c$ obtained from the gap and
$\hat\rho(\pi)$ are shown in figure~\ref{fltc}. One can see that
the difference between the two estimates shrinks and one can also
see that the average of the two estimates approach a finite limit
as the lattice spacing decreases.


The size and behavior of the
scaling violations we see are consistent with usual expectations,
reflecting contributions
of higher dimensional continuum operators to the action and the
operators. We leave for the future a more detailed investigation of
these scaling violations, but note that the orders of magnitude
are typical.



\begin{figure}[ht]
\centering
\includegraphics[width=\textwidth]{images/rho_diag.pdf}
\caption{ The square of the
$\hat\rho(\pi)$ as a function of $\ln(lt_c)$ along the crossing line. }
\label{rho_diag}
\end{figure}


\subsection{General features of $f_c (l)$.}

In summary, using all of our data, a rough estimate for the
critical size on our crossing line in the $(f,l)$ plane, is $l_ct_c\approx
0.6$. $l_c\approx \sqrt{\sigma}$, where $\sigma$ is the string
tension and we used the approximate relation $0.6 \sigma=t_c$. In
QCD units, $l_c$ is about half a fermi, in general accordance with the
accumulating body of numerical data which sets the transition from
perturbative to string-like behavior at rather short scales, of
less than one fermi. Clearly, varying the crossing line, the loop shape
and/or the details of the continuum construction of the Wilson
loop operator will produce slightly different values of $l_c$.
However, the small variation over scales in the critical line
$f_c( l)$ we find indicates that the transitions will occur in a
restricted range of scales. Numerically, it seems that $f_c(l)$ asymptotes
to a positive constant as $l\to 0$. This excludes choices of
Wilson loop definitions with $l_c t_c << 1$. However, if it is
somehow determined by
further simulations on very fine lattices that $f_c(l)\to 0$
as $l\to 0$, such choices must be recognized as possible;
in that case it is unlikely that an effective string description will
hold for $l$ too close to $l_c$. Even if
it were possible, we would not choose such a definition of regulated
continuum Wilson loop operators.

We have seen the scaling violations in the $F_c(L,b)$
approximations to $f_c (l)$ in a range of physical scales close to
the large $N$ transition point. It is impossible to decrease the
physical size of the loops much more, because that would require
too large
values of $b$ for the large $N$ twisted reduction trick
to work on a $7^4$ lattice. Similarly, it is impossible to get to
much larger loops because working at too small values of $b$ takes
the lattice system through a first order phase transition to a
strong coupling phase which has no relevance to continuum physics.

We found that
$f_c(l)$ flattened out for small loops. It is possible that
$f(l)$ does decrease to zero as $l$ goes to zero, but if this at
all happens, one would expect the effect to be logarithmic (i.e.
$f(l) \sim\frac{1}{|\log \Lambda l|}$), since $f(l)$ ought to
behave roughly as the gauge coupling constant at scale $l$, which
enters the perimeter divergence smearing is intended to eliminate.
When $l$ increases, $f_c(l)$ starts increasing more rapidly and
eventually, $f_c(l)$ would violate its lattice bound. It is a
strain to arrange for too large loops to have a spectral gap,
because for large loops confinement is reflected by an almost
uniform eigenvalue density.

Note however that confinement is not necessary in order to have
our large $N$ transitions; even if the eigenvalue density does not
approach a uniform distribution for very large loops it remains
possible that large loops have no gap in their spectrum. Thus, a
similar analysis might apply also to models that do not confine.
